%=============================================================
%  General Introduction
%=============================================================
\chapter*{General Introduction}
\addcontentsline{toc}{chapter}{General Introduction}

\section*{Context}

The Tunisian retail landscape is undergoing a profound digital transformation.
While e-commerce adoption has accelerated significantly since 2020, the market
remains largely fragmented: most local businesses still rely on informal social
commerce --- WhatsApp groups, Facebook marketplace posts, and Instagram direct
messages --- to reach their customers.
Those that do attempt a more structured online presence typically turn to
generic solutions such as WooCommerce or Shopify, which provide a catalogue
and a payment form but offer little in terms of intelligent personalization
or customer engagement.

In the cosmetics and personal care sector, this gap is particularly visible.
Customers face an overwhelming number of products with no guidance tailored
to their specific skin type, concerns, or purchasing history.
The result is decision fatigue, low conversion rates, and high return rates ---
problems that a well-designed, AI-augmented platform is uniquely positioned
to solve.

\section*{Problem Statement}

This project addresses the following question:
\textit{How can a full-stack e-commerce platform be designed to deliver both
a secure, reliable shopping experience and intelligent, personalized
interactions --- without relying on expensive third-party AI services or
compromising customer data privacy?}

Three concrete challenges arise from this question:
\begin{enumerate}
  \item \textbf{Architecture}: building a maintainable, scalable system using
        Clean Architecture principles that cleanly separates business logic
        from infrastructure concerns.
  \item \textbf{AI integration}: embedding a Retrieval-Augmented Generation
        (RAG) chatbot and a behavioral recommendation engine that operate
        entirely on local infrastructure, with no data sent to external providers.
  \item \textbf{Security}: implementing defense-in-depth --- from token revocation
        and brute-force protection to role-based access control and prompt
        injection guards --- without sacrificing developer experience.
\end{enumerate}

\section*{Proposed Solution}

\textbf{LUMINA} is a full-stack intelligent e-commerce platform built for the
Tunisian cosmetics market.
It combines a Next.js 14 frontend, an Express 5 / TypeScript backend following
Clean Architecture, and a Python/FastAPI RAG microservice running entirely
on local infrastructure (Ollama + sentence-transformers).

The platform delivers:
\begin{itemize}
  \item A complete shopping experience --- catalogue with advanced filtering,
        persistent cart, atomic checkout, and order tracking.
  \item Secure authentication --- email OTP verification, TOTP two-factor
        authentication, JWT blacklisting, and brute-force protection.
  \item An AI chatbot grounded in the product catalogue, protected by a
        three-layer prompt injection guard.
  \item A personalized recommendation engine driven by behavioral signals,
        requiring no external ML service.
  \item A full administration dashboard for product, user, and order management.
\end{itemize}

\section*{Report Structure}

This report is organized into seven chapters, following the three Scrum
iterations that structured the development process:

\begin{itemize}
  \item \textbf{Chapter 1 --- Project Context} presents the host company,
        a critical analysis of existing solutions, and the proposed project
        framed within the Scrum methodology.

  \item \textbf{Chapter 2 --- Requirements Specification} defines the
        functional and non-functional requirements, the system's actors,
        and the use case diagrams for each module.

  \item \textbf{Chapter 3 --- Conceptual Study} details the system architecture,
        Clean Architecture layers, sequence diagrams for key flows, and the
        domain class diagram.

  \item \textbf{Chapter 4 --- State of the Art} presents comparative studies
        for each technology choice --- frontend, backend, ORM, authentication,
        AI stack, and caching --- with justified selections.

  \item \textbf{Chapter 5 --- Iteration 1} covers the implementation of the
        e-commerce core: authentication, catalogue, cart, checkout, and
        basic administration.

  \item \textbf{Chapter 6 --- Iteration 2} details the AI and personalization
        layer: onboarding questionnaire, RAG chatbot with injection protection,
        and the behavioral recommendation engine.

  \item \textbf{Chapter 7 --- Iteration 3} describes the administration dashboard
        and the security hardening layer: JWT revocation, RBAC, and Zod
        input validation.
\end{itemize}
